
\subsection*{7.5. «~Conservation du nombre~».}
Supposons maintenant @@M0@@ propre, donc @@M1@@ et @@M2@@ propres sur @@M3@@ resp. @@M4@@, de sorte qu'on dispose d'homomorphismes 2.3 :
@@M5@@
Ceci dit, on a la relation
@@M6@@
(«~théorème de conservation du nombre~»), en d'autres termes, on a commutativité dans le diagramme
@@M7@@
Pour prouver (7.5.1), on est encore ramené au cas où @@M8@@, @@M9@@ sous-schéma fermé intègre de @@M10@@, et avec les notations de la démonstration de 7.4 on est ramené à prouver que
@@M11@@
ce qui est bien connu (EGA III 7.9.4, ou Exposé III 5.7).

\subsection*{7.6.}
Il est immédiat que l'on a
@@M12@@
ce qui entraîne, compte tenu que @@M13@@ est @@M14@@-linéaire, que le diagramme suivant est commutatif :
@@M15@@
